\section{Q-learning allocates capital across signals}\label{sec:14}

The outer learner makes a different decision: what fraction of one \$10 million budget should track each funded strategy? Its seven put/straddle targets are 100/0, 90/10, 75/25, 50/50, 25/75, 10/90 and 0/100. Weights sum to one; there is no outer leverage or cash action. A weekly decision based on prior-close information takes effect at the next close. The first session retains old holdings, then positions resize and drift until the next decision.


\subsection*{State, action and action-dependent holdings}

Nine state inputs describe four economic opportunities---put-spread rank strength, the put implied-versus-realized variance gap, straddle IV-versus-RV richness and completed straddle premium---plus relative strategy volatility, strategy-return correlation, VIX, relative recent Sharpe and previous straddle weight. Previous holdings matter because changing allocations incurs additional costs. A target is a fractional claim on a funded strategy index, not a new integer-contract order book. Underlying trading costs remain embedded in source returns; resizing adds source-based option and hedge charges.


\begin{equation*}R_w(\lambda)=\sum_{d\in w}(z_d-\lambda z_d^2),\qquad z_d=e_d/\widehat{\sigma}_{\mathrm{ref}}\end{equation*}

\subsection*{Reward and continuation value}

Reference volatility is estimated from the training window's fixed-50/50 excess returns. The parameter \(\lambda\) penalizes squared normalized fluctuations. At \(\lambda\) = 0, the directive is profit alone: additional normalized return earns reward without an explicit charge for the accompanying volatility. A higher-profit, higher-variance strategy can therefore attract capital even when its risk-adjusted performance is weaker. This objective does not maximize Sharpe. The discount \(\gamma\) = 0.99 separately weights continuation value; it is unchanged across risk preferences.


Monthly fitted Q iteration uses the preceding 104 completed weeks, with at least 52. A quadratic function of straddle weight has coefficients linear in standardized state, giving 30 coefficients. Each Bellman update regresses immediate reward plus discounted next-state maximum value onto the same basis \cite{ernst2005}. We use ridge coefficient 10 and 200 updates. Seven prior allocations crossed with seven actions supply counterfactual paths on each market week; they do not create independent market observations.


Q uses 50/50 before enough training weeks exist; each fixed control starts at its own target. A discount-zero learner supplies a matched one-period control. In the factorial experiment, Q is refitted independently for each entry policy. Strategy-dependent risk features, reward paths and resizing charges are rebuilt for that policy, while economic signal definitions remain shared.


This outer experiment addresses an extension proposed in the original paper: freeze each signal's proposals, then learn how to combine their exposures. Entry sizes change the inner funded portfolios; Q changes the capital assigned across them. The accounting layers are deliberately distinct, allowing us to remove either learning layer in the four-way comparison.


\FloatBarrier
